\documentclass[12pt]{article}
\usepackage{graphicx} % Required for inserting images
\usepackage{graphicx} % Required for inserting images
\usepackage[utf8]{inputenc}
\usepackage{amsmath, amssymb, amsthm}
\usepackage{graphicx,float}
\usepackage[T1]{fontenc}
\usepackage[thinc]{esdiff}
\usepackage{enumitem}

%Customization
\usepackage[letterpaper,top=0.2in,bottom=1in,left=0.8in,right=0.8in,heightrounded]{geometry}

%Line height
\renewcommand{\baselinestretch}{1.15}

%Parskip and Parindent
\setlength{\parindent}{0pt}
\setlength{\parskip}{0.7em}

\title{\textbf{Double Nested Operator Puzzle} \\[-0.6em]}
\author{Submission for Robols Math}
\date{}


\begin{document}
\maketitle
\vspace{-2em}

Solve for the expression:

\[ \prod_{k=67}^{6767} \frac{\pi}{2} \sum_{n=0}^{\infty} \frac{\frac {k}{2n+1}(-1)^n (\frac{k\pi}{2})^{2n} }{(2n)!} \]

\textbf{Solution:}

{Reorder the terms:}
\[ \prod_{k=67}^{6767} \sum_{n=0}^{\infty} \frac{\frac {1}{2n+1} \frac{k\pi}{2} (-1)^n (\frac{k\pi}{2})^{2n} }{(2n)!} \]
\[ \prod_{k=67}^{6767} \sum_{n=0}^{\infty} \frac{(-1)^n (\frac{k\pi}{2})^{2n+1} }{(2n+1)(2n)!} \]
\[ \prod_{k=67}^{6767} \sum_{n=0}^{\infty} \frac{(-1)^n (\frac{k\pi}{2})^{2n+1} }{(2n+1)!} \]

{Notice that the expression simplifies to the Taylor expansion of $\sin{\frac{k\pi}{2}}$. Then,}
\[ \prod_{k=67}^{6767} \sin{\frac{k\pi}{2}} \]
{Evaluate the expression}
\[\sin{\frac{(67)\pi}{2}} \cdot \sin{\frac{(68)\pi}{2}} \cdot \sin{\frac{(69)\pi}{2}} \cdot \cdots \cdot \sin{\frac{(6767)\pi}{2}}\]
{Since $\sin{\frac{68\pi}{2}}=\sin{34\pi}=0$, then the whole expression is equal to 0}.







\end{document}
